\documentclass[aspectratio=169,12pt]{beamer}

% ---------------------------------------------------------------------------
%  Step-by-step version of mu95_from_cms.tex: how the CMS-EXO-21-004 limits
%  become the mu95 used in the injection scan. One idea per slide.
%  Build: ../../build_slides.sh (figures, code boxes, then latexmk)
% ---------------------------------------------------------------------------

\usepackage{amsmath}
\usepackage{booktabs}
\usepackage{cancel}
\usepackage[T1]{fontenc}
\usepackage{lmodern}
\usepackage{listings}

% Code is shown straight from the source files (\lstinputlisting), never
% copied: snippets.tex is generated by code_snippets.py, which finds each
% line range by text. Run it before compiling.
\definecolor{codebg}{RGB}{247,247,247}
\definecolor{codekw}{RGB}{0,71,133}
\definecolor{codecom}{RGB}{110,110,110}
\definecolor{codestr}{RGB}{150,60,20}
\lstset{basicstyle=\ttfamily\scriptsize, backgroundcolor=\color{codebg},
  keywordstyle=\color{codekw}\bfseries, commentstyle=\color{codecom}\itshape,
  stringstyle=\color{codestr}, numbers=left, numberstyle=\tiny\color{softgray},
  numbersep=6pt, xleftmargin=2.2em, frame=single, rulecolor=\color{codebg},
  showstringspaces=false, columns=fullflexible, keepspaces=true,
  upquote=true, breaklines=false}
\lstdefinestyle{plainout}{basicstyle=\ttfamily\tiny, keywordstyle={},
  stringstyle={}, commentstyle={}}
\input{snippets.tex}
% where a snippet comes from, under the listing
\newcommand{\codesource}[1]{\par\vspace{0.2em}{\tiny\color{softgray} from #1}\par}
\graphicspath{{figs/}}

\usetheme{default}
\setbeamertemplate{navigation symbols}{}
\useinnertheme{circles}

\definecolor{cmsBlue}{RGB}{0,71,133}
\definecolor{accent}{RGB}{214,96,29}
\definecolor{softgray}{RGB}{120,120,120}
\definecolor{boxbg}{RGB}{238,243,249}

\setbeamercolor{frametitle}{fg=cmsBlue}
\setbeamercolor{title}{fg=cmsBlue}
\setbeamercolor{structure}{fg=cmsBlue}
\setbeamercolor{itemize item}{fg=accent}
\setbeamerfont{frametitle}{series=\bfseries}
\setbeamertemplate{frametitle}{%
  \vspace{0.6em}\insertframetitle\par
  {\color{accent}\rule{\textwidth}{1.2pt}}\par}
\setbeamertemplate{footline}{%
  \hfill{\color{softgray}\scriptsize\insertframenumber/\inserttotalframenumber}\hspace{1em}\vspace{0.5em}}

\newcommand{\muobs}{\mu_{95}^{\text{obs}}}
% the one-line takeaway at the bottom of a slide
\newcommand{\takeaway}[1]{\vfill
  \colorbox{boxbg}{\parbox{0.96\textwidth}{\centering\vspace{0.3em}#1\vspace{0.3em}}}
  \par\vspace{0.6em}}

\title{\textbf{What Signal Strength Is Still Allowed?}}
\subtitle{Turning the last CMS limit into a $\mu$ for the injection study, step by step}
\date{24 November 2026}

\begin{document}

% ---------------------------------------------------------------------------
\begin{frame}[plain]
  \titlepage
\end{frame}

% ===========================================================================
%  Part 1: the question
% ===========================================================================

\begin{frame}{Step 0: the question}
  \begin{itemize}
    \item The injection study asks: if signal hemispheres sneak into the mixing
          pool, how much do they change the mixed background?
    \item The answer depends on \textbf{how much signal} is there.
    \item So far we used the theory cross section. But CMS has already
          excluded that much signal.
  \end{itemize}
  \takeaway{We want the \textbf{largest amount of signal that data still allows}.}
\end{frame}

\begin{frame}{Step 1: what $\mu$ means}
  \begin{itemize}
    \item $\mu$ is the signal strength: how much signal, relative to the theory
          prediction.
          \[ \mu = \frac{\sigma_{\text{signal}}}{\sigma_{\text{theory}}} \]
    \item $\mu = 1$: exactly the theory cross section.
    \item $\mu = 0.01$: 1\% of it.
    \item $\mu = 0$: no signal at all.
  \end{itemize}
  \takeaway{$\mu$ is a ratio, so it does not care about units or energy.}
\end{frame}

\begin{frame}{Step 2: what a 95\% CL upper limit means}
  \begin{itemize}
    \item A search looks for a bump. If it sees none, it reports the largest
          signal that would still be consistent with what it saw.
    \item ``95\% CL upper limit'': any signal \emph{bigger} than this would have
          shown up at least 95\% of the time, so it is ruled out.
    \item Written as a signal strength, that limit is $\mu_{95}$:
      \begin{itemize}
        \item $\mu_{95} < 1$: the theory is excluded at this mass;
        \item $\mu_{95} > 1$: the theory is still allowed.
      \end{itemize}
  \end{itemize}
  \takeaway{$\mu_{95}$ is exactly the number we want: \textbf{the most signal
  data can still contain}.}
\end{frame}

% ===========================================================================
%  Part 2: where the numbers come from
% ===========================================================================

\begin{frame}{Step 3: which search}
  \begin{itemize}
    \item \textbf{CMS-EXO-21-004}: ``Searches for pair-produced multijet
          resonances using data scouting''.
    \item PRL 133 (2024) 201803, arXiv:2404.02992.
    \item 128\,fb$^{-1}$ of 13\,TeV scouting data: the Run 2 version of this
          analysis.
    \item It searched for gluino pairs, each decaying to 3 quarks (same signal
          as ours), and excluded gluinos lighter than 1.7\,TeV.
  \end{itemize}
  \takeaway{The most recent and strongest limit on our signal.}
\end{frame}

\begin{frame}{Step 4: where the numbers live}
  \begin{itemize}
    \item CMS papers put their plotted numbers on \textbf{HEPData}, a public
          database of tables.
    \item This paper: record \texttt{ins2774247}, table ``Figure 4c''
          (the three-jet-resonance limits).
    \item One row per gluino mass, 200 to 2000\,GeV.
    \item Saved locally:\\
          {\small\texttt{excluded\_xs/CMS-EXO-21-004\_HEPData\_Figure4c.csv}}
  \end{itemize}
  \takeaway{Everything that follows is arithmetic on this one table.}
\end{frame}

% ===========================================================================
%  Part 3: the arithmetic
% ===========================================================================

\begin{frame}{Step 5: the catch -- the limit is not on $\sigma$}
  \begin{itemize}
    \item The table does not limit the cross section $\sigma$ directly.
          It limits $\sigma \times B \times A$:
      \begin{itemize}
        \item $B$: branching ratio to 3 quarks. Here $B = 1$.
        \item $A$: \textbf{acceptance}, the fraction of signal events that pass
              the paper's selection.
      \end{itemize}
    \item Why? The data can only constrain events that were actually selected.
    \item Example, M300: $A \approx 0.001$. Only 1 gluino pair in 1000 passes the
          old selection.
  \end{itemize}
  \takeaway{The limit is on the \emph{visible} cross section, $\sigma \times A$.}
\end{frame}

\begin{frame}{Step 6: the trick -- the acceptance cancels}
  \begin{itemize}
    \item The table also gives the \textbf{theory} visible cross section,
          $\sigma_{\text{th}} \times A$, with the same $A$.
    \item Divide one by the other:
          \[ \mu_{95} \;=\; \frac{\text{limit on } \sigma A}{\sigma_{\text{th}} A}
             \;=\; \frac{\sigma_{95} \, \cancel{A}}{\sigma_{\text{th}}\, \cancel{A}}
             \;=\; \frac{\sigma_{95}}{\sigma_{\text{th}}} \]
    \item The acceptance drops out. We never need to know the old selection.
  \end{itemize}
  \takeaway{$\mu_{95}$ = (limit) $\div$ (theory), both read off the same row.}
\end{frame}

\begin{frame}{In the code: the division}
  Every limit is divided by the theory $\sigma A$ from the same row
  (last two lines):
  \vspace{0.3em}
  {\lstset{basicstyle=\ttfamily\tiny}\sniplimitdiv}
  \codesource{\srclimitdiv}
\end{frame}

\begin{frame}{Step 7: worked example, M300}
  \centering
  \begin{tabular}{lr}
    \toprule
    limit on $\sigma A$ (observed) & $1.49 \times 10^{-3}$ pb \\
    theory $\sigma A$             & $0.441$ pb \\
    \midrule
    $\mu_{95} = 1.49\times10^{-3} \,/\, 0.441$ & $\mathbf{3.39 \times 10^{-3}}$ \\
    \bottomrule
  \end{tabular}
  \vspace{1em}
  \begin{itemize}
    \item In events (multiply by 128\,fb$^{-1}$ = 128{,}000 pb$^{-1}$):
      \begin{itemize}
        \item theory predicts $\approx$ 56{,}000 selected signal events;
        \item data allows at most $\approx$ 190.
      \end{itemize}
  \end{itemize}
  \takeaway{At 300\,GeV only \textbf{0.34\%} of the theory signal can exist.}
\end{frame}

\begin{frame}{Step 8: worked example, M1000}
  \centering
  \begin{tabular}{lr}
    \toprule
    limit on $\sigma A$ (observed) & $3.74 \times 10^{-4}$ pb \\
    theory $\sigma A$             & $0.0125$ pb \\
    \midrule
    $\mu_{95} = 3.74\times10^{-4} \,/\, 0.0125$ & $\mathbf{0.030}$ \\
    \bottomrule
  \end{tabular}
  \vspace{1em}
  \begin{itemize}
    \item In events: theory predicts $\approx$ 1{,}600 selected signal events;
          data allows at most $\approx$ 48.
  \end{itemize}
  \takeaway{At 1\,TeV only \textbf{3\%} of the theory signal can exist.}
\end{frame}

\begin{frame}{In the code: the output for our two masses}
  Running \texttt{mu95\_table.py} prints one row per mass. Our two rows:
  \vspace{0.4em}
  {\lstset{basicstyle=\ttfamily\fontsize{4.6}{5.5}\selectfont}\sniplimitout}
  \codesource{\srclimitout}
  \takeaway{Read across: acceptance, theory $\sigma$ and $\sigma A$, the observed
  and expected limits; after the bar, $\mu_{95}$ observed, expected, and the
  expected band from $-2\sigma$ to $+2\sigma$.}
\end{frame}

\begin{frame}{Step 9: observed or expected?}
  \begin{itemize}
    \item The table gives two limits:
      \begin{itemize}
        \item \textbf{expected}: what the search would set if data had no
              signal, on average (with a $\pm1\sigma$, $\pm2\sigma$ range);
        \item \textbf{observed}: what it actually set with real data.
      \end{itemize}
    \item The expected limit describes sensitivity. The observed limit describes
          what the data allow.
    \item Ours: M300 observed $3.39\times10^{-3}$ vs expected $1.99\times10^{-3}$
          (inside $+1\sigma$); M1000 observed $0.030$ vs expected $0.038$.
  \end{itemize}
  \takeaway{We use the \textbf{observed} limit: it is what data actually allow.}
\end{frame}

\begin{frame}{Step 10: sanity check over all masses}
  \centering
  \includegraphics[height=0.56\textheight]{mu95_vs_mass.pdf}
  \takeaway{Same division at every mass. $\mu_{95}$ crosses 1 between 1.6 and
  1.8\,TeV, matching the paper's \textbf{1.7\,TeV} exclusion.}
\end{frame}

% ===========================================================================
%  Part 4: using it
% ===========================================================================

\begin{frame}{Step 11: from 13\,TeV to our 13.6\,TeV samples}
  \begin{itemize}
    \item The limit is at 13\,TeV; our MC is at 13.6\,TeV with a larger cross
          section.
    \item $\mu$ is a ratio, so we keep it and apply it to \emph{our} theory
          cross section:
          \[ \sigma_{95}(13.6) = \mu_{95} \times \sigma_{\text{th}}(13.6) \]
    \item M300: $0.00339 \times 637.6$ pb $=$ 2.16 pb.
    \item M1000: $0.030 \times 0.484$ pb $=$ 0.0145 pb.
  \end{itemize}
  \takeaway{Assumption: the allowed \emph{fraction} of the theory signal is the
  same at both energies.}
\end{frame}

\begin{frame}{Step 12: our own selection does not matter}
  \begin{itemize}
    \item Our samples are already weighted to $\sigma_{\text{th}}(13.6)$.
    \item Setting $\mu = \mu_{95}$ just multiplies every signal weight by
          $\mu_{95}$.
    \item Whatever cuts we apply afterwards (HT $>$ 550, mass asymmetry, \ldots)
          act on those scaled events, so our acceptance is automatically
          included.
  \end{itemize}
  \takeaway{The old analysis's acceptance cancels; ours comes for free.}
\end{frame}

% ===========================================================================
%  Part 4b: how every event is weighted
% ===========================================================================

\begin{frame}{Weights 1: why every MC event is divided by $N$}
  \begin{itemize}
    \item A simulated sample of $N$ events stands in for $\sigma \times L$ real
          events. So each simulated event counts as
          \[ w = \frac{\sigma \, L}{N}. \]
    \item Example: M300, $\sigma = 637.6$ pb, $N = 15{,}254$ events, so each
          event carries $637.6 / 15{,}254 = 0.0418$ pb.
    \item We fill with $L = 1$ pb$^{-1}$ (the analyzer default), so weights are in
          pb per event. Every plot here is unit-area, so $L$ drops out anyway.
  \end{itemize}
  \takeaway{Per dataset: $w_{\text{rel}} = \sigma / N$. The mixer calls this
  \texttt{w\_rel}.}
\end{frame}

\begin{frame}{Weights 2: what $N$ counts}
  \begin{itemize}
    \item $N$ must count \textbf{every} event that was generated, including the
          ones our cuts later throw away. Otherwise the selected events would be
          over-weighted.
    \item The slimmer records this as bin 0 of its \texttt{cutflow}
          histogram, ``All events'', before any cut:
  \end{itemize}
  \vspace{0.2em}
  {\lstset{basicstyle=\ttfamily\tiny}\snipsigcutflow}
  \codesource{\srcsigcutflow}
  $N$ = ``All events''; the $\geq$5-jet and HT cuts come after it.
\end{frame}

\begin{frame}{Weights 3: $N$ for the signal, and jet matching}
  \begin{itemize}
    \item The signal was made with MadGraph + Pythia and MLM jet matching. Matching
          rejects events that would be double counted; for M1000 it kept about
          20\% of the 50{,}000 generated events, 11{,}365 in all (last line):
  \end{itemize}
  \vspace{0.3em}
  \snipgenlog
  \codesource{\srcgenlog}
  \begin{itemize}
    \item The kept events are the sample, so $N = 11{,}365$, the same as
          ``All events'' on the last slide.
    \item $\sigma$ is the theory value (0.4842 pb), not MadGraph's LO matched
          0.17 pb. That is the standard normalisation.
  \end{itemize}
\end{frame}

\begin{frame}{In the code: reading $N$}
  The mixer reads bin 0 of each file's cutflow. A file without one is an error,
  not skipped: skipping it would shrink $N$ and inflate every weight.
  \vspace{0.2em}
  {\lstset{basicstyle=\ttfamily\tiny}\snipreadn}
  \codesource{\srcreadn}
\end{frame}

\begin{frame}{In the code: building \texttt{w\_rel}}
  Those counts are summed over every file of a dataset (\texttt{n0}); the last
  line is $\sigma / N$:
  \vspace{0.2em}
  {\lstset{basicstyle=\ttfamily\tiny}\snipwrel}
  \codesource{\srcwrel}
\end{frame}

\begin{frame}{Weights 4: the actual numbers}
  \snipwtable
  \codesource{\srcwtable}
  \begin{itemize}
    \item \textbf{x4:} the four copies of the signal file are all counted in $N$,
          so each signal event carries 1/4 of the x1 weight.
  \end{itemize}
\end{frame}

\begin{frame}{Weights 5: a mixed event has two parents}
  \begin{itemize}
    \item A mixed event is built from two hemispheres, the ``seed'' and the
          ``match'', which can come from different datasets.
    \item The mixer stores the \textbf{product} of the two parents' weights:
  \end{itemize}
  \vspace{0.3em}
  \snipxsweight
  \codesource{\srcxsweight}
  \takeaway{\texttt{xs\_weight} $= w_{\text{rel}}(\text{seed}) \times
  w_{\text{rel}}(\text{match})$, one value per mixed event.}
\end{frame}

\begin{frame}{Weights 6: what the analyzer fills with}
  \begin{itemize}
    \item A product of two $\sigma/N$ has units pb$^2$. The analyzer takes the
          square root (\texttt{weight\_power} = 0.5) to get back to pb per event:
  \end{itemize}
  \vspace{0.3em}
  \snipfillweight
  \codesource{\srcfillweight}
  \takeaway{Fill weight $= \sqrt{w_{\text{rel}}(\text{seed}) \times
  w_{\text{rel}}(\text{match})}$, the geometric mean of the two parents.}
\end{frame}

\begin{frame}{Weights 7: QQ, SQ and SS go into separate histograms}
  Count how many parents come from a signal file (ID $\geq 5482$), then fill the
  QQ, SQ or SS histogram with the event's own weight. Keeping them apart is what
  lets $\mu$ be applied afterwards, without rerunning:
  \vspace{0.2em}
  {\lstset{basicstyle=\ttfamily\tiny}\snippairsplit}
  \codesource{\srcpairsplit}
\end{frame}

\begin{frame}{Weights 8: real numbers for single events}
  \snipwpairs
  \codesource{\srcwpairs}
  \begin{itemize}
    \item Signal strength $\mu$ multiplies the signal parent's $w_{\text{rel}}$.
          The fill weight is a square root, so an SQ event changes by
          $\sqrt{\mu}$: at $\muobs$, $0.01447 \to 0.00084$ (a factor
          $\sqrt{0.00339} = 0.058$).
    \item An SS event has two signal parents, so it changes by $\mu$.
  \end{itemize}
  \takeaway{That is exactly the $\mu^{1/2}$ and $\mu$ on the next slide.}
\end{frame}

\begin{frame}{In the code: how $\mu$ scales the signal}
  The same scaling, applied to the finished QQ, SQ and SS histograms:
  \vspace{0.4em}
  \snipcombine
  \codesource{\srccombine}
  With $\sqrt{w}$ filling \texttt{p} = 0.5, so this is
  $\text{QQ} + \mu^{1/2}\,\text{SQ} + \mu\,\text{SS}$.
  \takeaway{The cuts were applied when the histograms were filled, so our
  acceptance is already inside \texttt{parts}.}
\end{frame}

\begin{frame}{Step 13: what we scanned}
  \centering
  \begin{tabular}{lcc}
    \toprule
    curve & M300 & M1000 \\
    \midrule
    $\muobs$                 & 0.0034 & 0.030 \\
    $10 \times \muobs$       & 0.034  & 0.30 \\
    $100 \times \muobs$      & 0.34   & 3.0 \\
    theory ($\mu = 1$)       & 1      & 1 \\
    \bottomrule
  \end{tabular}
  \vspace{1em}
  \begin{itemize}
    \item The x4 studies duplicate each signal hemisphere 4 times, so they apply
          $\mu/4$ to match x1.
  \end{itemize}
  \takeaway{$10\times$ and $100\times$ show how far we are from any effect.}
\end{frame}

\begin{frame}{In the code: choosing the $\mu$ values}
  \texttt{-{}-mu95} turns one number into the list of curves:
  $\mu_{95}$, $10\,\mu_{95}$, $100\,\mu_{95}$, plus 0 and 1.
  \vspace{0.4em}
  \snipmuargs
  \codesource{\srcmuargs}
  \takeaway{One option, one number per mass; everything else follows.}
\end{frame}

\begin{frame}{In the code: running all four studies}
  The run script gives each mass its own $\mu_{95}$, and each x4 study its
  factor 1/4 (\texttt{-{}-mu-scale 0.25}):
  \vspace{0.4em}
  \sniprunloop
  \codesource{\srcrunloop}
  \takeaway{One command reproduces every plot: \texttt{zsh run\_mu95\_scan.sh}.}
\end{frame}

\begin{frame}{In the code: $\Delta\chi^2$}
  Per $\mu$: shape vs $\mu = 0$ bin by bin, minus the signal MC's own noise.
  \vspace{0.2em}
  {\lstset{basicstyle=\ttfamily\tiny}\snipdchisq}
  \codesource{\srcdchisq}
\end{frame}

\begin{frame}{In the code: the printed result, M300 x1}
  \vspace{0.4em}
  \snipscanout
  \codesource{\srcscanout}
  \vspace{0.4em}
  Columns: $\mu$, signal share of the weight, $\Delta\chi^2$/bins,
  $\sqrt{\Delta\chi^2}$ (known-shape significance), generic-test significance.
  \takeaway{At $\mu_{95}$ = 0.00339: $\Delta\chi^2$ = 0.2 (HT) and 0.0 (mass).
  2$\sigma$ needs 4.}
\end{frame}

\begin{frame}{Step 14: the result}
  \begin{itemize}
    \item At $\muobs$ the change to the mixed background is
          $\Delta\chi^2 \le 0.3$ for every observable.
    \item Seeing a $2\sigma$ effect needs $\Delta\chi^2 = 4$
          (see the $\chi^2$ explainer).
    \item M300 would need 60--100$\times$ its allowed cross section to reach
          $2\sigma$ in mass; M1000 never gets there in the scan.
  \end{itemize}
  \takeaway{At the signal levels data still allow, signal in the mixing pool
  \textbf{does not visibly change} the background.}
\end{frame}

\begin{frame}{Summary: the recipe}
  \begin{enumerate}
    \item Find the latest limit: CMS-EXO-21-004, HEPData ``Figure 4c''.
    \item Read two numbers from the row for your mass: the observed limit on
          $\sigma A$ and the theory $\sigma A$.
    \item Divide: $\mu_{95}$ = limit / theory. The acceptance cancels.
    \item Check: $\mu_{95} = 1$ near 1.7\,TeV, as the paper says.
    \item Multiply our signal weights by $\mu_{95}$ and run the scan.
  \end{enumerate}
  \vspace{0.5em}
  \centering
  M300: $\muobs = 3.39\times10^{-3}$ \qquad M1000: $\muobs = 0.030$
\end{frame}

\begin{frame}{Reproducing everything}
  \scriptsize
  \begin{tabular}{@{}p{0.5\textwidth}p{0.46\textwidth}@{}}
    \toprule
    file & what it does \\
    \midrule
    \texttt{excluded\_xs/mu95\_table.py} & HEPData CSV $\to$ $\mu_{95}$ at every mass \\
    \texttt{run\_mu95\_scan.sh} & the table, all four scans, the zoom, merged PDFs \\
    \texttt{../plot\_injection\_scan.py} & the scan
      (\texttt{-{}-mu95}, \texttt{-{}-xlim}) \\
    \texttt{slides/mu95\_from\_cms/make\_figures.py} & the $\mu_{95}$-vs-mass plot \\
    \texttt{weights/weights\_table.py} & per-dataset and per-event weights, from
      the mixer's file tables (copied from cmslpc to the work directory) \\
    \texttt{weights/signal\_cutflow.py} & the signal files' cutflows (run on cmslpc) \\
    \texttt{slides/mu95\_from\_cms/code\_snippets.py} & pulls the code into these slides \\
    \texttt{build\_slides.sh} & rebuilds the figures, the code boxes and the decks \\
    \bottomrule
  \end{tabular}
  \vspace{1em}

  \small Paths relative to \texttt{run3-mj-analyzer/scripts/signal\_injection/}
  (see its \texttt{README.md}); plots and histograms live in
  \texttt{\textasciitilde/analysis/histograms/inject\_26\_11\_23/}. Every code box
  in this deck is read from those files when the deck is compiled; the line
  numbers are the files' own.
\end{frame}

\end{document}
